\section{Lista de características: seguimiento de tareas estructurado}

La lista de características es uno de los componentes más críticos de la solución completa; resuelve directamente el problema en que los Agentes declaran prematuramente la finalización e intentan ejecutar todas las tareas de una sola vez.

\subsection{Principio de diseño}

El Agente Inicializador genera, basado en el prompt inicial del usuario, una lista exhaustiva de requisitos funcionales. En el ejemplo de "clonar claude.ai", dicha lista incluye más de 200 puntos funcionales, por ejemplo: "el usuario puede abrir un nuevo chat, ingresar una pregunta, presionar Enter y ver la respuesta de la IA".

Cada punto funcional se marca inicialmente como \texttt{"passes": false}, de modo que el posterior Agente de Codificación pueda observar claramente el panorama completo de las funciones implementadas, sin llegar a creer erróneamente que el proyecto está terminado simplemente porque algunas funciones ya han sido ejecutadas.

\subsection{Selección del formato JSON}

\begin{lstlisting}[language={}, caption={Estructura JSON de un punto funcional individual dentro de la lista de características}]
{
  "category": "functional",
  "description": "El botón de Nuevo chat crea una conversación fresca",
  "steps": [
    "Navegar a la interfaz principal",
    "Hacer clic en el botón 'Nuevo chat'",
    "Verificar que se haya creado una nueva conversación",
    "Comprobar que el área de chat muestre el estado de bienvenida",
    "Verificar que la conversación aparezca en la barra lateral"
  ],
  "passes": false
}
\end{lstlisting}

\begin{importantbox}{¿Por qué elegir JSON en lugar de Markdown?}
El equipo de Anthropic descubrió tras experimentar que, comparado con archivos Markdown, los modelos tienen menos probabilidades de modificar o sobrescribir inapropiadamente el contenido de archivos JSON. Las características estructuradas del JSON hacen que el Agente tienda a modificar únicamente el valor del campo \texttt{passes}, en lugar de reescribir toda la descripción. Esta es una técnica valiosa de ingeniería de prompts: \textbf{la elección del formato de datos influye en las tendencias conductuales del modelo}.
\end{importantbox}

\subsection{Reglas estrictas de modificación}

Al Agente de Codificación se le instruye explícitamente que solo pueda modificar el valor del campo \texttt{passes} —de \texttt{false} a \texttt{true}. El equipo de Anthropic utilizó un lenguaje contundente en el prompt:

\begin{quote}
\textit{"No es aceptable eliminar o editar las pruebas porque esto podría llevar a funcionalidad faltante o con errores."}
\end{quote}

Este énfasis en la redacción resulta sumamente efectivo para controlar el comportamiento del Agente. Le indica claramente al Agente que eliminar o modificar las descripciones de funciones es inaceptable, previniendo fundamentalmente que el Agente adopte atajos para "completar" tareas bajando los estándares.

\begin{warningbox}{La tendencia al "engaño" del Agente}
Sin restricciones explícitas, un Agente podría "completar" una tarea más rápidamente eliminando o reduciendo los estándares de las pruebas. Esto es análogo a la Ley de Goodhart en humanos: cuando una métrica se convierte en el objetivo mismo, deja de ser una buena métrica. En el diseño de sistemas de Agentes, es necesario prevenir este comportamiento mediante reglas estrictas y criterios de evaluación inmutables.
\end{warningbox}

\subsection{Resumen de la sección}

La lista de características materializa los requisitos del proyecto en puntos funcionales rastreables mediante un formato JSON estructurado; cada función cuenta con pasos de verificación claros. La elección del formato JSON, el estado inicial marcado completamente como fallido y las reglas estrictas que prohíben modificar las descripciones de funciones conforman conjuntamente un mecanismo eficaz para evitar que los Agentes declaren prematuramente la finalización del trabajo.

